% !TEX root = main.tex

\section{Literature Review}
\label{sec:literature}

\subsection{Classical time-series models and the limits of return predictability}

The autoregressive integrated moving average framework of \citet{box1970time}
remains the canonical linear model for the conditional mean of a univariate
series and the natural benchmark against which richer specifications are judged.
Its known limitation in financial applications lies in the second moment: daily
returns exhibit conditional heteroskedasticity, which motivated the ARCH model of
\citet{engle1982autoregressive} and its GARCH generalisation
\citep{bollerslev1986generalized}, and which implies that a correctly specified
linear conditional-mean model may still misstate predictive uncertainty by a wide
margin. A second limitation concerns the conditional mean itself. In a
comprehensive out-of-sample audit, \citet{welch2008comprehensive} found that most
proposed return predictors fail to improve upon the historical average, and
\citet{timmermann2008elusive} characterised return predictability as local and
episodic---present in short-lived pockets rather than as a stable population
property. Forecast combination offers partial insurance against this instability
\citep{rapach2010out}, but the collective evidence establishes a demanding
baseline: any gain claimed for a more elaborate model must survive genuine
out-of-sample evaluation over heterogeneous regimes.

\subsection{Machine learning approaches}

Machine learning methods relax the linearity restriction at the cost of a larger
effective parameter space. Random forests \citep{breiman2001random} control
overfitting through bootstrap aggregation of decorrelated trees and provide
internal measures of predictor importance, properties that make them a common
nonlinear benchmark in return forecasting. In the most systematic comparison to
date, \citet{gu2020empirical} showed that regression trees and neural networks
yield economically significant improvements in measuring equity risk premia,
attributing the gains chiefly to nonlinear predictor interactions that linear
methods miss. For sequence models, the LSTM architecture
\citep{hochreiter1997long} was shown by \citet{fischer2018deep} to outperform
random forests and logistic regression in classifying daily S\&P~500 constituent
returns, although the documented profitability of the strategy declined in later
subperiods---an instability consistent with the pocketed predictability described
above. Subsequent applied work has extended these architectures with
macroeconomic inputs and hybrid convolutional--recurrent designs
\citep{latif2025integrating}. Two limitations recur throughout this literature:
the fitted models are difficult to interpret, and their outputs are point
predictions unaccompanied by a distributional statement.

A third limitation is less often discussed and is central to this paper. Because
these models are typically reported at a single regularisation setting chosen by
the author rather than selected on held-out data, it is difficult to distinguish
a model that has found signal from one that has fitted noise. Section
\ref{sec:mean} shows that the distinction is decisive here.

\subsection{Bayesian inference and uncertainty quantification}

The Bayesian literature addresses the distributional limitation directly.
\citet{west2020bayesian} surveys multivariate Bayesian forecasting with emphasis
on scalability and structure uncertainty, arguing that decision-relevant
forecasts require the full predictive distribution rather than a summary of its
centre. \citet{qiu2018multivariate} develop the multivariate Bayesian structural
time series model, in which local linear trends, seasonal components and
spike-and-slab selection over a candidate pool of predictors jointly capture
contemporaneous correlation across target series while guarding against
overfitting. In the return-predictability setting specifically, Bayesian
treatments of model uncertainty---averaging over predictor subsets rather than
conditioning on a single specification---were shown by \citet{avramov2002stock}
and \citet{cremers2002stock} to improve out-of-sample performance precisely
because no individual model is reliably correct. Closer to the present study,
\citet{chandra2021bayesian} apply Bayesian neural networks to multi-step stock
price forecasting and report usable predictive intervals even through the
COVID-19 volatility spike, while \citet{wong2024quantifying} show that standard
uncertainty-quantification schemes for neural networks miscalibrate under
volatility clustering, underscoring that calibration must be verified rather than
assumed. Proper scoring rules \citep{gneiting2007strictly} and tests of equal
predictive accuracy \citep{diebold1995comparing} supply the formal apparatus for
such verification, and this paper applies both throughout.

For coefficient shrinkage we adopt the horseshoe prior of
\citet{carvalho2010horseshoe}, the continuous analogue of the spike-and-slab
selection used by \citet{qiu2018multivariate}, sampled by the auxiliary-variable
scheme of \citet{makalic2016simple}.

\subsection{Distribution shift and the fragility of fixed-origin evaluation}

A strand of the machine learning literature that has had limited uptake in
empirical finance concerns dataset shift: the case in which the joint
distribution of predictors and target differs between estimation and deployment
\citep{quinonero2009dataset}. The concern is directly relevant here, because the
macro-financial variables most commonly used to condition equity forecasts---an
index level, a policy rate, a yield---are precisely the series whose means move
permanently across a monetary regime change. Standardising such a series on a
fixed estimation window and applying the transformation to a later window
produces inputs far outside the range the estimator has seen, and any model with
an unbounded output head extrapolates along them. Section
\ref{sec:pitfall} documents the consequences and the repair.

\subsection{Hybrid architectures and the remaining gap}

A growing strand seeks to combine the representational capacity of recurrent
networks with the distributional discipline of econometric volatility models.
\citet{kim2018forecasting} integrate LSTM networks with multiple GARCH-type
specifications and report improved index volatility forecasts relative to either
component alone, and \citet{wang2022garlic} document analogous gains for a
combined LSTM--GARCH-family model in a highly volatile price series. These
hybrids, however, typically attach a volatility filter to a point-forecasting
network rather than deliver a coherent posterior predictive distribution, and
published comparisons rarely hold the information set fixed across model
families.

The present study is positioned in that gap. It evaluates classical, machine
learning and Bayesian models on a single panel with an identical macro-financial
predictor set and an identical regime-spanning out-of-sample window, and it
scores them on both point accuracy and the calibration of their stated
uncertainty. To the extent that the literature reviewed above establishes weak
and unstable conditional-mean predictability alongside strong and persistent
structure in conditional variances, the comparison speaks to a practical
question left open by prior work: whether the principal contribution of
macroeconomic drivers and Bayesian inference at the daily horizon lies in sharper
point forecasts or in more honest predictive distributions. The decomposition in
Section~\ref{sec:decomposition} answers it quantitatively.
